\begin{table}[htbp]
\centering
\small
\begin{tabular}{lccccc}
\hline\hline
\textbf{Modelo} & \textbf{N} & \textbf{Prec} & \textbf{Rec} & \textbf{F1} & \textbf{MSE} \\
\hline
SNN (A--B) & 100 & 0.160 & 0.472 & 0.239 & 0.790 \\
SNN (A--B) & 200 & 0.104 & 0.299 & 0.155 & 0.859 \\
SNN (A--B) & 400 & 0.083 & 0.239 & 0.124 & 0.892 \\
\hline
SNN (A--B--C) & 100 & 0.263 & 1.000 & 0.416 & 0.737 \\
SNN (A--B--C) & 200 & 0.264 & 1.000 & 0.417 & 0.734 \\
SNN (A--B--C) & 400 & 0.264 & 0.994 & 0.417 & 0.730 \\
\hline\hline
\end{tabular}
\caption{Resultados de detección de anomalías en dataset CALLT2. N representa el número de neuronas en las capas B y C para los modelos SNN. Se muestran los mejores resultados obtenidos tras optimización con Optuna para cada configuración.}
\label{tab:resultados-callt2-escalabilidad}
\end{table}